\section{Evaluación y puntos de referencia}

\subsection{Sistema de métricas de evaluación}

\begin{itemize}
    \item \textbf{Sample efficiency}: cuántas transitions se consumen en promedio por cada update;
    \item \textbf{Reward saturation}: si el plateau de la curva de reward es estable;
    \item \textbf{Entropy decay}: una caída rápida de la entropy indica un colapso de la exploration;
    \item \textbf{Policy divergence}: si la KL y la dirección del gradiente se mantienen consistentes con la critic signal.
\end{itemize}

\begin{importantbox}{Definir con cuidado la ventana de evaluación}
Al probar cada checkpoint hay que usar una semilla aleatoria fija y un número fijo de rollout steps, para evitar que un single-run high reward induzca a conclusiones erróneas. Se recomienda reportar la puntuación de best-of-$N$ junto con la media y la desviación estándar.
\end{importantbox}

\subsection{Ejemplos de Benchmark}

\begin{itemize}
    \item \textbf{MuJoCo locomotion}: PPO como baseline, SAC ofrece mayor eficiencia de muestreo;
    \item \textbf{Atari}: A3C/A2C mostraron que el paralelismo por lotes puede acortar significativamente la convergencia;
    \item \textbf{Robotics sim2real}: se controlan estrictamente la KL, el safety filter y la entropy;
    \item \textbf{StarCraft multiagente}: centralized critic + decentralized actor.
\end{itemize}

\begin{knowledgebox}{Errores frecuentes al hacer benchmark}
Reproducir el best score usando una sola single random seed puede ocultar el policy collapse y la training instability. Se recomienda ejecutar al menos entre 3 y 5 semillas, y graficar juntas la trayectoria de la KL y la curva de reward.
\end{knowledgebox}

\subsection{Resumen de la sección}

La evaluación no debe limitarse a observar el final reward, sino también monitorear la entropy, la KL, la critic loss y el sampling ratio, para poder determinar si el Actor-Critic es realmente estable.

